\documentclass[twoside]{article}

\usepackage{ustj}
\usepackage{pdflscape}
\usepackage{multicol}
\usepackage{caption}

\addbibresource{mss.bib}

\newcommand{\authorname}{Jack Fox}
\newcommand{\authorpatp}{\patp{nomryg-nilref}}
\newcommand{\affiliation}{FoxyLabs}

%  Make first page footer:
\fancypagestyle{firststyle}{%
\fancyhf{}% Clear header/footer
\fancyhead{}
\fancyfoot[L]{{\footnotesize
              %% We toggle between these:
              Manuscript submitted for review.\\
              % {\it Urbit Systems Technical Journal} I:1 (2024):  1–3. \\
              ~ \\
              Address author correspondence to \authorpatp.
              }}
}
%  Arrange subsequent pages:
\fancyhf{}
\fancyhead[LE]{{\urbitfont Urbit Systems Technical Journal}}
\fancyhead[RO]{Obelisk}
\fancyfoot[LE,RO]{\thepage}

%%MANUSCRIPT
\title{Obelisk:  Reinventing SQL \\ for Modern Computing}
\author{\authorname~\authorpatp \\ \affiliation}
\date{}

\begin{document}

\maketitle
\thispagestyle{firststyle}

\begin{abstract}
Obelisk is a novel relational database management system (\textsc{rdbms}) and its scripting language, urQL, designed for the experimental Urbit computing platform but adaptable to other systems. Building on E.\ F.\ Codd’s relational algebra and unencumbered by backward compatibility constraints, Obelisk is a temporal relational database derived from first principles. Key innovations include time travel for querying historical database states, proper set implementation, elimination of nulls, idempotent queries, and enforced composability. Obelisk ensures all user operations are accessible via both the urQL scripting language and \textsc{api}, reflecting Urbit’s pure functional architecture. Obelisk’s time-travel feature indexes every schema and content state change by time, ensuring idempotent queries and atomic script execution. This paper explores how Obelisk advances relational database design while addressing \textsc{sql}’s deficiencies, offering a modern, principled approach to structured data management.
\end{abstract}

% We will adjust page numbering in final editing.
\pagenumbering{arabic}
\setcounter{page}{1}

\tableofcontents

\section{Introduction}

The Obelisk \textsc{rdbms} and urQL scripting language\footnote{We prefer \emph{scripting language} over \emph{query language} for semantic reasons. It does more than query.} run on the experimental Urbit computer and network but could just as well be implemented on any other computer---except for certain features dependent on Hoon language features or the Urbit network architecture.\footnote{Notably subnet security features and shrubbery enabling \lstinline[style=inlinecode]{SELECT} statements.} Like all relational databases, it is built upon E.\ F.\ Codd's pioneering computer science work developing the relational algebra and uniting it with set theory and predicate calculus.\footnote{\citet{Codd1970}; \citet{Chamberlin1974} Codd's seminal paper focuses primarily on studying set relations and does not invent a query language, but sets out the parameters for an effective and efficient language. Predicate calculus was already the obvious way to filter stored data. \citet{Chamberlin1974} treat it as a given, discussing their syntax for applying it.} As relational databases implementing \textsc{sql} came to dominate structured data storage the \textsc{sql} standard evolved, always bearing, and hindered by, the necessary constraint of backward compatibility. As the first and only \textsc{rdbms} for a new computing paradigm, Obelisk has no obligation to the past and so is free to build anew from first principles with all the learnings of the last half century while perpetuating the best features of \textsc{sql}.

Obelisk's urQL, while clearly a dialect of \textsc{sql}, fixes numerous deficiencies in \textsc{sql}\footnote{E.g. the \textsc{sql} standard does not truly implement sets by default. Obelisk does.}, elevates composability to a syntactic requirement---not merely an option, adds one completely new major feature, time travel,\footnote{\citet{DarwenDate1998}; \citet{DarwenDate2002}. While influential in other aspects of Obelisk's design. We are less familiar with their work on temporal databases, so their research on temporal databases was not especially influential. Without any direct knowledge of his work, we credit the reputation of Rich Hickey.} and leaves no functionality inaccessible to the normal user.

The \textsc{rdbms} evolution that is Obelisk was greatly inspired by the Urbit architecture. Urbit is a computer whose entire state (OS and user space) is a pure function of its prior state.\footnote{\citet{Whitepaper}. Is any sufficiently complex computer function completely pure? Urbit and Obelisk both rely on crash to reject invalid commands that parsing and the \textsc{api} type system cannot otherwise validate.} The Urbit computer itself is an \textsc{acid} database. The entire state of the server is one noun\footnote{The \emph{noun} is an Urbit abstraction that maps beautifully to English. It is either a natural number, or an ordered pair of nouns. Thus the fundamental data structure of Urbit (and Obelisk) is a binary tree of natural numbers.} that advances only through pure functions. urQL's syntax is more than a ``better" scripting language, it is an accurate reflection of the underlying database server's operations. Anything the user can accomplish through urQL they can also do through the \textsc{api}\footnote{Of the \textsc{sql} \textsc{rdbms}s that do have \textsc{api}s, they are uniformly bolted-on \mbox{afterthoughts}.}, and vice versa. Together with Obelisk's native time travel capabilities and idempotence, there are no mysteries confronting the user, no functionality outside the scope of the \textsc{api}/scripting language like performing audits with database logs. There are no logs. (There are system views that emulate logs to support the familiar log analysis pattern.) Obelisk's state is both its current state and all prior states.\footnote{With the sole exception of dropping a database from the server, in which case all traces of the database vanish from the server.}

Obelisk's major innovations include:

\begin{itemize}
\item Time travel
\item Proper sets
\item No nulls
\item Composability
\end{itemize}

\section{Time Travel}

\begin{quote}
The totality of data in a data bank may be viewed as a collection of time-varying relations \citep{Codd1970}.
\end{quote}

In Obelisk time is primary and fundamental. Every change of state, whether to a database's schema or content, is indexed by time. Thus every query is idempotent because each query is implicitly or explicitly associated with a particular state in the series. And time is a series of steps, like in physical reality. A time step may be appended to the series, but there is no inserting, updating, or deleting time steps.

The rules enforcing time primacy in the Obelisk database engine are simple. Each database holds a most recent schema\footnote{Currently the only schema altering commands available are \lstinline[style=inlinecode]{CREATE} and \lstinline[style=inlinecode]{DROP}. \lstinline[style=inlinecode]{ALTER} is on the roadmap for an undetermined future release.} time (by user data table) and a most recent content time (by user data table). Every subsequent state change, whether to schema or content must be subsequent to the latest of the two times. Normally the user never needs to concern themselves with this requirement. When not overridden by \lstinline[style=inlinecode]{AS OF} the database engine applies the host ship's time\footnote{A \emph{ship} is the generic term for an Urbit computer on the Urbit network, which has a unique, immutable, cryptographically protected identity. The ship's time corresponds to the computer time on the Urbit VM's host computer, whether Linux/Unix, Mac, or Windows.} to determine the current state of the database.

Script execution, regardless of the number or sequencing of commands, potentially alters the state of the database as if everything happened \emph{all at once}, even though the commands execute sequentially. Users only need  to be aware of this rule when applying \lstinline[style=inlinecode]{AS OF} to override \lstinline[style=inlinecode]{NOW}. It is possible for a script to result in different ending state times for different tables.  Violation of time constraints (or any other error) causes the entire script to fail. (Scripts are always atomic.)

The second and final rule is once a script introduces a query---a \lstinline[style=inlinecode]{SELECT} command returning data results (but not any intermediate \lstinline[style=inlinecode]{SELECT})---all subsequent commands must also be queries. Further schema or data changes are not allowed.

Among the metadata returned by queries are the schema and content times (labelled \lstinline[style=inlinecode]{schema time} and \lstinline[style=inlinecode]{data time}) used by the engine to create the query results. A query executes with a \emph{de facto} time of the latest of the two for the state of each underlying table or view. In order to recreate the same query one needs to specify this \lstinline[style=inlinecode]{AS OF} time. Thus every query is implicitly or explicitly idempotent. 

Creating a new schema component, \lstinline[style=inlinecode]{DATABASE}, \lstinline[style=inlinecode]{NAMESPACE},\footnote{A \emph{namespace} is known as \emph{schema} in the rest of the \textsc{sql} world, leading to confusion.} \lstinline[style=inlinecode]{TABLE}, or \lstinline[style=inlinecode]{VIEW}, may be backdated (or forward dated) to any time subsequent to the latest of schema time or content time (\lstinline[style=inlinecode]{CREATE DATABASE} is of course entirely independent). Users override schema and data manipulation commands at their own peril. It is possible to effectively lock all or portions of the database for update.

\begin{lstlisting}[style=listingcode_sql]
CREATE DATABASE db2 AS OF ~2099.1.1 ;

%obelisk-result:
  %results
    [ %message 'created database %db2' ]
    [ %server-time ~2024.9.29..21.14.00..7fc8 ]
    [ %schema-time ~2099.1.1 ]
\end{lstlisting}

\noindent
Getting passing results for all operations on the database now requires a time subsequent to this future creation date.

\begin{lstlisting}[style=listingcode_sql]
(crash)
CREATE NAMESPACE: %ns1 as-of schema time out of order
\end{lstlisting}

And the database is not even visible in the current context.

\begin{lstlisting}[style=listingcode_sql]
FROM sys.sys.databases AS OF NOW
SELECT * ;

%obelisk-result:
  %results
    [ %message 'SELECT' ]
    %result-set
      database  sys-agent  sys-tmsp  data-ship
       ↳ data-agent  data-tmsp
      db1  /gall/dojo  ~2024.9.26
       ↳ ~zod  /gall/dojo  ~2024.9.26
      sys  /gall/dojo  ~2024.9.26
       ↳ ~zod  /gall/dojo  ~2024.9.26
    [ %server-time ~2024.9.26..22.03.39..2233 ]
    [ %message 'sys.sys.databases' ]
    [ %schema-time ~2024.9.26 ]
    [ %data-time ~2024.9.26 ]
    [ %vector-count 2 ]
\end{lstlisting}

We can see it in a future time context.

\begin{lstlisting}[style=listingcode_sql]
FROM sys.sys.databases AS OF ~2099.1.1
SELECT * ;

database  sys-agent  sys-tmsp  data-ship
 ↳ data-agent  data-tmsp
animal-shelter  /gall/dojo  ~2024.10.1
 ↳ ~zod  /gall/dojo  ~2024.10.1
db1  /gall/dojo  ~2024.9.26
 ↳ ~zod  /gall/dojo  ~2024.9.26
db2  /gall/dojo  ~2099.1.1
 ↳ ~zod  /gall/dojo  ~2030.1.1
sys  /gall/dojo  ~2024.9.26
 ↳ ~zod  /gall/dojo  ~2024.9.26
\end{lstlisting}

Manipulating \lstinline[style=inlinecode]{AS OF} time for schema components is rarely useful, but could be used for creating and populating a database with history. What is much more useful is setting the content time for query or content manipulation\footnote{Setting the \lstinline[style=inlinecode]{AS OF} time when manipulating data (\lstinline[style=inlinecode]{INSERT}, \lstinline[style=inlinecode]{UPDATE}, \lstinline[style=inlinecode]{DELETE}) results in a new content state based on the state in effect at that time.} relative to some past state, or mixing times to audit content changes.

\begin{lstlisting}[style=listingcode_sql]
FROM my-table T1
OUTER JOIN my-table AS OF <some past date> T2
WHERE NOT EXISTS T1.key
   OR NOT EXISTS T2.key
SELECT * ;
\end{lstlisting}

\noindent
(The example demonstrates urQL’s non-standard join and predicate syntax. A full explanation follows.)

The astute reader may have already inferred the last missing piece of the time travel puzzle. When specified in a \lstinline[style=inlinecode]{FROM} clause, the database engine will use the closest content state less than or equal to the time requested, using the table schema in a state likewise less than or equal to that time.

Say \lstinline[style=inlinecode]{my-table} has content mutated at times \lstinline[style=inlinecode]{~2024.10.2.. 16.54.41..df15} and \lstinline[style=inlinecode]{~2024.10.4..17.32.11..a1c2}.

\begin{lstlisting}[style=listingcode_sql]
FROM db2..my-table-1 AS OF ~2024.10.3
SELECT * ;

%obelisk-result:
  %results
    [ %message 'SELECT' ]
    %result-set

      <results not shown>

    [ %server-time ~2024.10.30..17.00.48..85c2 ]
    [ %message 'db1.dbo.my-table-1' ]
    [ %schema-time ~2023.1.4..14.15.21..ad17 ]
    [ %data-time ~2024.10.2..16.54.41..df15 ]
    [ %vector-count 3 ]
\end{lstlisting}

\noindent
As explained above, the database engine returns metadata for this query which allows us to identify the implicit idempotent query. Any \lstinline[style=inlinecode]{AS OF} time \lstinline[style=inlinecode]{>= ~2024.10.2..16.54.41..df15} and \lstinline[style=inlinecode]{< ~2024.10.4..17.32.11..a1c2} is the same idempotent query.

\begin{lstlisting}[style=listingcode_sql]
FROM db2..my-table-1 AS OF ~2024.10.2..16.54.41..df15
SELECT * ;
\end{lstlisting}

\noindent
Identifying \emph{the} idempotent query with the one exactly matching the content state time is important. The current last content state in the table's series can always be superseded by a content state one tick later.

Idempotent queries have a big technical benefit. The engine can cache query results with never a thought given to cache expiration. Simple systems have fewer bugs and less attack surface. Obelisk does this for views---which in Obelisk are read only. Making a long-running query a view requires only a one-time performance hit...until the backing table data state progresses, in which case there is another one-time performance hit.

\section{The Relational Model}

\subsection{Relations}

So far we have casually referred  to \emph{tables} and \emph{views} as standard components of relational databases. In this section we introduce Codd's term, \emph{relation}.\footnote{\citet{Codd1970}.} (What follows is a hasty summary of Codd's formal definition, and it is somewhat more restrictive than Codd's definition in order to make a point further below.) Consider a proper set of data tuples (i.e. every member unique). All the tuples have the same length, and positionally each element across the tuples is of the same type. This is a \emph{relation}. Arranging the tuples as rows in a 2D matrix (arbitrarily ordered), we can index the columns by unique labels.\footnote{Here we deviate from Codd's formal definition, where he only suggests the columns---he has formally defined them as domains---may be labeled, the domains are not necessarily unique. One must bear in mind Codd's work is theoretic. He even points out that he is not designing a data query language. He is providing the foundation upon which to design query languages.} We can refer to each element in any of the tuples by the column label and we have our definition of \emph{relation} for the discussion below.

\subsection{Proper sets}

There is no \lstinline[style=inlinecode]{DISTINCT} keyword. All relations in Obelisk are proper sets, as Codd intended.

A \emph{table} in Obelisk is a set of immutable persistently stored relations---indexed by time. Rows are indexed by a unique key referencing one or more column names, which ensures a proper set. Mutating the contents results in a new subsequent relation.

A \emph{view} is a persistently \emph{defined} relation. It is defined by \lstinline[style=inlinecode]{SELECT} commands and set operators over tables and views. The relation's set-like characteristic is maintained by shrinking the rows to a proper set. Views are indexed by time. A change to the definition or state of any table referenced in the definition results in a new time index. As mentioned earlier, if a user queries a view the system caches its state for that point in time. We distinguish between \emph{persistently defined} and \emph{persistently stored} because there is no requirement for view data to persist. Due to query idempotence it can always be recreated.

The \lstinline[style=inlinecode]{SELECT} command itself is a special case. It defines a relation, but it is also a presentation specification, and the resulting presentation may not be a relation (by our restricted definition), even though the resulting rows always form a proper set. Here is an example.

\begin{lstlisting}[style=listingcode_sql]
FROM my-table
SELECT column-1, column-1, column-1
\end{lstlisting}

\noindent
This command defines a \emph{proper} relation consisting of the column named \emph{column-1}, but the presentation layer repeats it three times. Thus the resulting presentation is not a (restricted) relation because there are three columns all with the same label.

This \lstinline[style=inlinecode]{SELECT} does result in a (restricted) relation

\begin{lstlisting}[style=listingcode_sql]
FROM my-table
SELECT column-1, column-1 AS Col-2, column-1 AS Col-3
\end{lstlisting}

\noindent
This quibbling is for the sake of emphasizing (our) formal definition of \emph{relation}, and to point out when defining a view the \lstinline[style=inlinecode]{SELECT} command must provide for unique column names. Informally one may consider the results of \lstinline[style=inlinecode]{SELECT} as a \emph{pseudo-relation}. 

Trivially \lstinline[style=inlinecode]{SELECT} without a \lstinline[style=inlinecode]{FROM} clause is a one-row relation because Obelisk always assigns unique labels to literals if not otherwise aliased.

\begin{lstlisting}[style=listingcode_sql]
SELECT 0, ~nomryg-nilref, ~2025.9.24

%obelisk-result:
  %results
    [ %message 'SELECT' ]
    %result-set
      literal-0  literal-1  literal-2
      0  ~nomryg-nilref  ~2025.9.24
    [ %server-time ~2025.9.24..17.56.02..de98 ]
    [ %schema-time ~2025.9.22..22.03.28..6f07 ]
    [ %data-time ~2025.9.22..22.03.28..6f07 ]
    [ %vector-count 1 ]
\end{lstlisting}

\noindent
(When there is no \lstinline[style=inlinecode]{FROM} clause schema and data times reference the creation time of the system database.)

\subsection{Joins}

The urQL join\footnote{We respect the purity of Codd's original relational algebra, thus there is no \emph{inner join} only join. Outer joins are a set union on a joined relation as well as the remaining rows not being joined (\citeauthor{Codd1990}, \citeyear{Codd1990}, p.~107). Nonetheless outer joins have great practical benefit and we opt to dispense with superfluous keywords. Codd never published original research on outer joins, but after their introduction by other authors \citep{Chamberlin1974}, he covered them in a book \citep{Codd1990}. Oracle is widely credited with the first adoption of outer joins in the early 1980s, and they were first included in the \textsc{ansi} \textsc{sql}-92 standard. Outer joins are not included in the first Obelisk beta.} is restricted to the key words \lstinline[style=inlinecode]{JOIN}, \lstinline[style=inlinecode]{CROSS JOIN}, \lstinline[style=inlinecode]{LEFT JOIN}, \lstinline[style=inlinecode]{RIGHT JOIN}, \lstinline[style=inlinecode]{OUTER JOIN}.  \lstinline[style=inlinecode]{JOIN} is either a natural join or a $\theta$-join.\footnote{\citet{Codd1970}.}

The absence of \lstinline[style=inlinecode]{ON <predicate>} syntactically determines a natural join, which in Obelisk must adhere to the following rules.  

Natural join is always between two relations, the subject of the current \lstinline[style=inlinecode]{JOIN} clause and any relation previously specified in the \lstinline[style=inlinecode]{FROM} clause, so order of relations matters. The joining characteristic is equality between all like-named and like-typed columns common between the two.\footnote{\citet{Codd1970} describes this as a cross join filtered for equality on the operands.} Hence a natural join may result in a relation of more rows than either of the two input relations. If there are matching columns with more than one previously listed relation, the join fails and a $\theta$-join is required.

The availability of natural joins encourages transparent and composable schema designs, but bear in mind common metadata columns such as \emph{timestamp} must have uniquely qualified names across tables and views in order for natural joins to be useful.

\sloppy
Syntactically speaking, a $\theta$-join is a join with \lstinline[style=inlinecode]{ON <predicate>} clause wherein at least one equality or inequality operator exists with a column from the current relation on one side and a column from a previously defined relation on the other, and such an operation must exist on both sides of all \lstinline[style=inlinecode]{OR} conjunctions.

The outer joins \lstinline[style=inlinecode]{LEFT}, \lstinline[style=inlinecode]{RIGHT}, and \lstinline[style=inlinecode]{OUTER} (full outer join) are substantially the same as for any other \textsc{sql} implementation and may be natural or $\theta$-joins.  \lstinline[style=inlinecode]{CROSS JOIN} is the cartesian join.

\subsection{No Nulls}

To avoid the quagmire of three-valued (or even four-valued) logic, nulls are entirely banned from urQL.\footnote{\citeauthor{DarwenDate1998} (1998; 2002)
% \citet{DarwenDate2002}
influenced this decision.} To begin with, columns may not be defined as nullable. This leaves outer joins as the other way traditional \textsc{sql} database engines can produce nulls. Once again \lstinline[style=inlinecode]{SELECT} is a special case and does not always produce \emph{ad hoc} relations\footnote{It is beyond the scope of this paper to invent a new formal definition of \emph{relation}. We imagine that instead of a relation consisting of one tuple type of ordered atomic values, there would be a root tuple type of ordered unique label, type, value triples. (Type ensures strong typing. An untyped version is more general.) Any number of unique sub-types consisting of any subset of these triples in the same order may also be present. Arranged vertically all the tuples align on label/type matches. Any resulting 2D matrix having more than the root tuple type present has some number of \emph{holes} in every occurrence of a sub-type, and is thus a \emph{sparse relation}.} in \textsc{sql} or urQL.

The purpose of a data management system is the organization and retrieval of data. The presentation of that data is an ancillary concern. As a \emph{from first principles} reinvention of relational data management, Obelisk is not obligated to return query results artificially shoehorned into a two-dimensional matrix.

Rather than queries returning column rows of fixed length, Obelisk returns vectors of unique label, type, and atomic datum. The labels---usually the column name from the underlying relation qualified by relation to distinguish like-named columns---and order of the triples are derived from the \lstinline[style=inlinecode]{SELECT} command. This extends to set operations, notably \lstinline[style=inlinecode]{UNION}.

Outer joins produce a list of variable length vectors. Let us call each potential variable length a \emph{vector type}. All column name/type labeled values potentially exist in one vector type, and smaller types have a proper subset of these labels in the same order. The upper bound for types in a \lstinline[style=inlinecode]{LEFT JOIN} or \lstinline[style=inlinecode]{RIGHT JOIN} is two and for \lstinline[style=inlinecode]{OUTER JOIN} is three. In the case of set unions the relation types may share no column types whatsoever---in which case there is a hypothetical root type never instantiated.

Presentation software is free to interpret a list of vectors in any manner. By vertically aligning labels, presentation layers are free to insert \lstinline[style=inlinecode]{NULL} or any other token in holes.

Returning to a previous query example

\begin{lstlisting}[style=listingcode_sql]
FROM my-table T1
OUTER JOIN my-table AS OF <some past date> T2
WHERE NOT EXISTS T1.key
   OR NOT EXISTS T2.key
SELECT * ;
\end{lstlisting}

\noindent
returns the following vector list in pseudo-hoon

\begin{lstlisting}[style=listingcode_sql]
[%vector [p=%t1-key q=[p=~.ud q=1]] [p=%t1-col-2 ...]
 ↳ ...<remaining columns>...]
[%vector [p=%t2-key q=[p=~.ud q=2]] [p=%t2-col-2 ...]
 ↳ ...<remaining columns>...]
...
\end{lstlisting}

\noindent
and in the following section we explain the remaining non-standard attribute of this query.

\subsection{Predicate Calculus}

To deal with column type sparseness resulting from outer joins and relation unions, urQL introduces two predicate keywords:

\begin{itemize}
\item  \lstinline[style=inlinecode]{EXISTS} This is NOT the standard \textsc{sql} \lstinline[style=inlinecode]{EXISTS} keyword. Rather it is a unary operator that tests for existence of a column in a possibly sparse relation row, or it can test that a \textsc{cte} returns data.

\item  \lstinline[style=inlinecode]{EQUIV} This binary operator is the same as equality, except comparison of two non-existent columns and/or \textsc{cte}s results in \lstinline[style=inlinecode]{TRUE}.
\end{itemize}

\noindent
The standard equality and inequality operators comparing any non-existing column and/or \textsc{cte} results in \emph{false}.

\subsection{Relational Integrity}

We plan to implement relational integrity by means of foreign key indices in a future release, substantially following the \textsc{sql} Standard.

\section{Composability}

\begin{quote}
The universality of the data sublanguage lies in its descriptive ability \citep{Codd1970}.
\end{quote}

The first rule of Composability Club is ``have a query syntax that looks compositional.''

\begin{lstlisting}[style=listingcode_sql]
WITH <ctes>
FROM
<joins>
SCALARS
WHERE
GROUP BY
HAVING
SELECT
ORDER BY
TOP
\end{lstlisting}

\noindent
This is the logical flow of data and operations defined by the \textsc{sql} standard (with a notable Obelisk addition discussed below) It is how users and the database engine agree on constructing proper results. Actual query execution plans may deviate in some way from rigidly following this logical path for performance reasons. This is why, unlike standard \textsc{sql}, urQL syntax sets the \lstinline[style=inlinecode]{SELECT} clause in its logical place, not before the \lstinline[style=inlinecode]{FROM} clause.

\subsection{Common Table Expressions}

Common Table Expressions entered the \textsc{sql} standard\footnote{The \textsc{ansi}/\textsc{iso} \textsc{sql} standard defined Common Table Expressions in 1999.} to make it easier for users to reason about sub-queries. Inlining subqueries makes for bad readability. The \textsc{sql} Standard must maintain inlining backwards compatibility. urQL has no such constraint. It enforces \textsc{cte} composability by banning all subquery inlining.

\sloppy
urQL requires a \textsc{cte} in all cases involving subqueries through runtime validation. Any usage in a predicate or scalar function that can accept a column value or literal also accepts a \textsc{cte} alias provided that the \textsc{cte} returns a singleton, one row by one column. This is not the case in the \lstinline[style=inlinecode]{SELECT} clause. Selecting a singleton \textsc{cte} is functionally equivalent to a \lstinline[style=inlinecode]{CROSS JOIN}, which urQL syntax prefers in the interest of explicit programming. Scalar function and predicate operators consuming a list (set) of values require the \textsc{cte} to return a single column of data. 

\begin{lstlisting}[style=listingcode_sql]
WITH Foo AS (...)
...
WHERE bar IN Foo
\end{lstlisting}

\subsection{Scalar Functions}

Scalar functions present another opportunity to enforce composability. The \lstinline[style=inlinecode]{SCALARS} keyword introduces a clause where all scalar functions must be defined. They are subsequently referenced by alias. A scalar function operates on any data in a single data row as well as singleton and list \textsc{cte}s and returns one consistently typed value.

\begin{lstlisting}[style=listingcode_sql]
SCALARS
  Sc1 IF foo = 1 THEN bar ELSE foobar END-IF
  Sc2 IF Sc1 = 'what-not' THEN %.y ELSE %.n END-IF
  Sc3 IF Sc2 THEN dt-col ELSE ~2025.9.6 END-IF
\end{lstlisting}

\noindent
In legacy \textsc{sql} this would all be nested and inlined in a \lstinline[style=inlinecode]{SELECT} clause or predicate. In urQL one inlines the scalar alias.

Obelisk supports \lstinline[style=inlinecode]{IF}, \lstinline[style=inlinecode]{CASE}, \lstinline[style=inlinecode]{COALESCE}, and many other common \textsc{sql} scalar functions.

\section{Security}

The Urbit architecture provides for single sign-on (\textsc{sso}) through the immutable identity of all computers on the network. Obelisk security\footnote{Not in the first beta release. The first beta prevents all access by a foreign ship.} allows for grouping identities, much like security on other \textsc{sql} database implementations. In addition parent, child, and sibling are predefined groups providing for subnet security. The Urbit network architecture provides for subnets with two levels of hierarchy---parent, a ship of identity type \emph{galaxy}, \emph{star}, or \emph{planet}, and children, ships of identity type \emph{moon}.

User permissions are for \emph{ships} (computers) on the network, for \emph{agents} (userspace programs), or a combination of the two.\footnote{Userspace security is still a work in progress in the Urbit architecture.} Permissions can be any combination of whitelists or blacklists affecting databases, namespaces, tables and views, or columns. Permissions may be read-only or for write affecting data. Schema read is allowed, but update can only be performed by the host ship. Whitelisting and blacklisting can be tailored by timeframe, both in regards to the data or schema state permitted or restricted and the effective time range of the permission itself.

Obelisk, like other \textsc{sql}-based relational databases, is susceptible to \textsc{sql} injection attacks, and the mitigation for Obelisk is the same as for other \textsc{dbms}s---implementing stored procedures and adding a level of SP permissions.

Short of dropping a database, any write---even dropping a table---results in a new state. Referencing a prior state and writing to it (with proper planning) can effectively correct any destructive write.

Obelisk security is believed to be robust up to disk and \textsc{ram} hacking.

\section{The Shrubbery}

Key to Codd's invention of the relational database is his critique of hierarchical databases.\footnote{\citet{Codd1970}, ``1.1.3 Access Path Dependence''. Codd uses the term \emph{data store} which is mostly interchangeable with \emph{database}, even if they conjure up different images in the mind.} It took another 20 years to mostly fulfill the demise of accessing data via paths. Yet the concept lingers and is not without merit. Urbit being both a computer fundamentally supporting a hierarchical (tree) architecture and a network of nodes with immutable identities is the perfect substrate for implementing a truly global namespace.\footnote{Not in the first beta release. The shrubbery concept of a universal namespace is an organic outgrowth of the Urbit file system (Clay). Within Urbit circles the concept is still amorphous. It could be more than an Urbit network \textsc{uri}, with endpoints functioning as data or software. It is hard to identify a single author or inventor. \patp{rovnys-ricfer} originated the term in conversation with \patp{littel-wolfur} and was the first to write about it (\patp{rovnys-ricfer}, personal communication, \textasciitilde 2025.10.18; \citeauthor{Blackman_YoungShrub}, undated). \patp{sarpen-laplux} developed the earliest working implementation in 2023. \patp{hastuc-dibtux} subsequently carried out extensive prototyping and advocacy (\citeyear{Fitzgerald_Shrub}). \patp{migrev-dolseg}'s \texttt{\%hawk} gall app is the most readily available and extensive implementation in this tradition to date (\citeyear{HanlenHawk}).}

Obelisk fully supports global namespace access to data and schema. This urQL/Obelisk functionality is not easily duplicated on computers other than Urbit.

\begin{lstlisting}[style=listingcode_sql]
SELECT DATABASE animal-shelter
\end{lstlisting}

\noindent
\lstinline[style=inlinecode]{SELECT DATABASE} returns the entire \emph{animal-shelter}\footnote{\citet{Levin2020}. Available for several database platforms, this is a readily available non-trivial sample database that demonstrates natural primary table keys.} database noun (see above).\footnote{Backups may be implemented as a copy of the database noun on another ship. The Urbit computer is already an \textsc{acid} database and nouns are de-duplicated on the same ship. So references to the same noun are pointers to the single physical occurrence of that noun.}

Formatted as a path in the global Urbit namespace (beginning with ship identity)

\begin{lstlisting}[style=listingcode_sql]
~nomryg-nilref/%obelisk/animal-shelter
\end{lstlisting}

\noindent
any caller on the network with sufficient security permissions can retrieve a copy of the database.

Working our way down the schema hierarchy we have access to namespaces within a database,

\begin{lstlisting}[style=listingcode_sql]
SELECT NAMESPACE animal-shelter.reference

~/%obelisk/animal-shelter/reference
\end{lstlisting}

\noindent
and pause to point out the path is on the local ship, but still confined by local desk\footnote{A \emph{desk} can be thought of as an application space inside of user space.} permissions. The returned noun is all the user defined tables, views, and indices in the selected namespace, including all their schema and data states.

\begin{lstlisting}[style=listingcode_sql]
SELECT my-table

~/%obelisk/db1/dbo/my-table
\end{lstlisting}

\noindent
The reader may wonder about the syntax here. Submission of urQL scripts to the Obelisk parser requires the name of a default database along with the script, and \lstinline[style=inlinecode]{dbo} is always the default namespace. The returned noun is the same metadata and result set whose print format we have seen above.

We cannot tell from the syntax (or the returned metadata) whether the query returned the data from a table or view. Views can shadow tables. To force a table add the \lstinline[style=inlinecode]{TABLE} keyword to the script, or in the namespace,

\begin{lstlisting}[style=listingcode_sql]
~/%obelisk/db1/dbo/table/my-table
\end{lstlisting}

\noindent
and use \lstinline[style=inlinecode]{AS OF} for a particular state,

\begin{lstlisting}[style=listingcode_sql]
~/%obelisk/db1/dbo/my-table/~2020.1.26
\end{lstlisting}

The final drill down in the namespace is to the column level,

\begin{lstlisting}[style=listingcode_sql]
SELECT db1..my-table.column-4
~/%obelisk/db1/dbo/my-table/column-4
\end{lstlisting}
\noindent
Keep in mind Obelisk returns proper sets of data rows. Only a view is capable of deterministic ordering of the result rows when accessing data in this manner.

Shrubbery support provides for hybrid relational/hierarchical data processing. The Hawk Urbit implementation of namespace computing provides for scripting language or \textsc{api} usage of Obelisk as well as path access and post processing. There is already a working adaptation of a \textsc{sql} Studio-like interface to Obelisk available in Hawk. 

\section{Data Structures}

\begin{quote}
Bad programmers worry about the code. Good programmers worry about data structures and their relationships (Torvalds, 2016).
\end{quote}

To better understand the workings of Obelisk we turn to the data structures. Obelisk is a work in progress approaching a beta release, so some structures still may change over time.

\subsection[\textsc{ast}/\textsc{api}]{AST/API}

The Obelisk parser parses urQL commands into the structures in file \lstinline[style=inlinecode]{sur/ast.hoon}, which also serves as the Obelisk \textsc{api}. We show the most important schema and content manipulating commands in Figure~\ref{fig:obelisk_ast_schema}. A \lstinline[style=inlinecode]{qualified-column} is one qualified by a \lstinline[style=inlinecode]{qualified-} \lstinline[style=inlinecode]{table}.

\begin{landscape}
\begin{figure}[p]
\centering
\vspace*{-0.8cm} % Adjust for header space
\begin{lstlisting}[style=listingblock, basicstyle=\setmonofont\footnotesize,numbers=none, xleftmargin=3em, multicols=3]
:: create-database
$:  %create-database
    name=@tas
	  as-of=(unit as-of)
==

:: drop-database         
$:  %drop-database
    name=@tas
    force=?
==

:: create-namespace
$:  %create-namespace
    database-name=@tas
    name=@tas
    as-of=(unit as-of)
==

:: create-table
$:  %create-table
    table=qualified-table
    columns=(list column)
    pri-indx=(list ordered-column)
    foreign-keys=(list foreign-key)
    as-of=(unit as-of)
==

:: predicate
(tree predicate-component)
:: qualified-table
$:  %qualified-table
    ship=(unit @p)
    database=@tas
    namespace=@tas
    name=@tas
    alias=(unit @t)
==

:: drop-table
$:  %drop-table
    table=qualified-table
    force=?
    as-of=(unit as-of) 
==

:: update
$:  %update
    ctes=(list cte)
    table=qualified-table
    as-of=(unit as-of)
    $:  columns=(list qualified-column)
        values=(list value-or-default)
    ==
    =predicate
==




:: delete
$:  %delete
    ctes=(list cte)
    table=qualified-table
    as-of=(unit as-of)
    =predicate 
==

:: truncate-table
$:  %truncate-table
    table=qualified-table
    as-of=(unit as-of)
==

:: predicate-component
 $? ops-and-conjs
	 qualified-column
	 unqualified-column
	 dime
	 value-literals
    cte-alias
    scalar-alias
    aggregate 
==

:: column
$:  %column
    name=@tas
    type=@ta
==
\end{lstlisting}
% \vspace*{-0.5cm} % Adjust for footer space
\caption{Obelisk schema creation and modification \textsc{ast}/\textsc{api} structures, from \texttt{/sur/ast}.}
\label{fig:obelisk_ast_schema}
\end{figure}
\end{landscape}

Obelisk does not yet support interacting with tables on foreign ships, so qualification is currently only by database and namespace. An alias may be assigned in the urQL. There is no distinction in urQL or the \textsc{ast} between tables and views. A view may shadow a table in the context of queries, but schema and data manipulating commands always operate on the specified table.

We see that \lstinline[style=inlinecode]{DELETE} requires a predicate (unlike standard \textsc{sql}) and the predicate data structure is tree-shaped. Predicate nodes are operators, conjunctions, or any number of representations of data.

Figure~\ref{fig:obelisk_query_ast} presents the \lstinline[style=inlinecode]{selection} data structure, also a tree structure which accommodates either a set operator or one of three other structures as nodes. \lstinline[style=inlinecode]{insert} and \lstinline[style=inlinecode]{merge} are for future implementation using the \lstinline[style=inlinecode]{%into} operator. The \lstinline[style=inlinecode]{query} data structure is central to reading and formatting persisted data and can be combined with any of the regular set operators (\lstinline[style=inlinecode]{%union}, \lstinline[style=inlinecode]{%except},\lstinline[style=inlinecode]{%intersect},\lstinline[style=inlinecode]{%divided-by}). If there are \textsc{cte}s, they are processed first to return one relation each. Prior \textsc{cte} relations are accessible by subsequent \textsc{cte}s. Internally different kinds of relation data structure represent data returned by \textsc{cte}s (\lstinline[style=inlinecode]{cte-alias}) or more commonly \lstinline[style=inlinecode]{qualified-table}, which includes (possibly table-shadowing) views. (\lstinline[style=inlinecode]{query-row} is a placeholder for future temporary table urQL syntax.)

\begin{landscape}
\begin{figure}[p]
\centering
\vspace*{-0.8cm} % Adjust for header space
\begin{lstlisting}[style=listingblock, basicstyle=\setmonofont\footnotesize,numbers=none, xleftmargin=3em, multicols=3]
:: selection
$:  %selection
    ctes=(list cte)
    set-functions=(tree set-function)
==

:: set-function
$?(set-cmd set-op) 

:: query
$:  %query
    from=(unit from)
    scalars=(list scalar)
    =predicate
    group-by=(list grouping-column)
    having=predicate
    =select
    order-by=(list ordering-column) 
==

:: set-op
$?  %union
    %except
    %intersect
    %divided-by
    %divide-with-remainder
    %into
==


:: set-cmd
$%(insert merge query)  

:: cte
$:  %cte
    name=@tas
    =query 
==

:: from
$:  %from
    =relation
    as-of=(unit as-of)
    joins=(list joined-object)  
==

:: scalar
$:  %scalar
    scalar=scalar-function
    alias=@t
==

:: select
$:  %select
    top=(unit @ud)
    columns=(list selected-column)  
==



:: selected-column
$%  qualified-column
    unqualified-column
    selected-aggregate
    selected-value
    selected-all
    selected-all-object 
==

:: insert
$:  %insert
	  table=qualified-table
	  as-of=(unit as-of)
	  columns=(unit (list @tas))
	  values=insert-values
==

:: relation
$%  qualified-table
    cte-alias
    query-row
==

:: joined-object
$:  %joined-object
    join=join-type
    =relation
    as-of=(unit as-of)
    =predicate
==
\end{lstlisting}
% \vspace*{-0.5cm} % Adjust for footer space
\caption{Obelisk query and selection \textsc{ast}/\textsc{api} structures, from \texttt{/sur/ast}.}
\label{fig:obelisk_query_ast}
\end{figure}
\end{landscape}

\lstinline[style=inlinecode]{selected-column} deals with all the different ways columns and data may be selected in urQL, qualified and unqualified columns, literals, all columns across all joins, and all columns of a specified relation. (Aggregate functions are for a future release.)

A potential enhancement to the urQL syntax and \textsc{ast} in a future release is to support zero-ary structures in the \lstinline[style=inlinecode]{SELECT} clause of a query. This is an area of research enabled by Obelisk's ability to produce result vectors and sub-vectors (see outer joins above). The model is how the Datalog language defines zero-ary structures, beginning a sub-structure with \lstinline[style=inlinecode]{:-} colhep and ending it with a \lstinline[style=inlinecode]{.} dot, allowing nested sub-structures.

The example

\begin{lstlisting}[style=listingcode_sql]
SELECT column-1, column-2, :- column-3, column-4., column-5
\end{lstlisting}
\noindent
would return one \lstinline[style=inlinecode]{vector} of \lstinline[style=inlinecode]{column-1}, \lstinline[style=inlinecode]{column-2}, \lstinline[style=inlinecode]{column-5}, and one \lstinline[style=inlinecode]{vector} of \lstinline[style=inlinecode]{column-3}, \lstinline[style=inlinecode]{column-4} for every instance of the same \lstinline[style=inlinecode]{column-1}, \lstinline[style=inlinecode]{column-2}, \lstinline[style=inlinecode]{column-5}. This would open the door to returning complex nouns from queries.

\subsection{Persistence Layer}

At the very top of the Persistence Layer (Figure\ref{fig:obelisk_persistence_ast}) is the server data structure.

\begin{lstlisting}[style=listingblock, label={lst:obelisk_persistence}]
+$  server  (map @tas database)
\end{lstlisting}

\noindent
This is simply an index into all the databases, including the \lstinline[style=inlinecode]{%sys} database tracking other databases. This is the noun which the Obelisk app persists. Alternate implementations using an app structure other than the one provided with Obelisk merely need to persist this noun.

\begin{landscape}
\begin{figure}[p]
\centering
\vspace*{-0.8cm} % Adjust for header space
\begin{lstlisting}[style=listingblock, basicstyle=\setmonofont\footnotesize,numbers=none, xleftmargin=3em, multicols=3]
:: database
$:  %database
    name=@tas
    created-provenance=path
    created-tmsp=@da
    sys=((mop @da schema) gth)
    content=((mop @da dataa) gth)
    =view-cache  
==

:: view-cache
((mop ns-obj-key cache) ns-obj-comp)  

:: ns-obj-key
$:  ns=@tas
    obj=@tas
    time=@da
==

:: cache
$:  %cache
    tmsp=@da
    content=(unit cache-content)
==


:: ns-obj-comp
|=  [p=data-obj-key q=data-obj-key]
^-  ?

:: cache-content
$:  rowcount=@
    rows=(list (map @tas @))
==

:: schema
$:  %schema
    provenance=path
    tmsp=@da
    =namespaces
    tables=(map [@tas @tas] table)
    views=((mop data-obj-key view)
      ns-obj-comp)
==

:: data
$:  %data
    ship=@p
    provenance=path
    tmsp=@da
    files=(map [@tas @tas] file)
==
:: table
$:  %table
    provenance=path
    tmsp=@da
    =column-lookup
    type-lookup=(map @tas @ta)
    pri-indx=index
    columns=(list column)
    indices=(list index)  
==

:: column-lookup
(map @tas [aura @])

:: view
$:  %view
    provenance=path
    tmsp=@da
    =selection
    =column-lookup
    type-lookup=(map @tas @ta)
    columns=(list column)
    ordering=(list column-order)  
==
\end{lstlisting}
% \vspace*{-0.5cm} % Adjust for footer space
\caption{Obelisk persistence layer \textsc{ast}/\textsc{api} structures, from \texttt{/sur/server-state}.  (Continues on next page.)}
\label{fig:obelisk_persistence_ast}
\end{figure}
\end{landscape}

\begin{figure}[h]
\ContinuedFloat
\centering
% \vspace*{-0.8cm} % Adjust for header space
\begin{lstlisting}[style=listingblock, basicstyle=\setmonofont\footnotesize,numbers=none, xleftmargin=3em, multicols=2, keepspaces=true, columns=flexible]
:: file
$:  %file
    ship=@p
    provenance=path
    tmsp=@da
    =column-addrs
    rowcount=@
    pri-idx=(tree [(list @) (map @tas @)])
    indexed-rows=(list indexed-row)
    =column-catalog  
==
:: indexed-row
$:  %indexed-row
    key=(list @)
    data=(map @tas @)
==





::
\end{lstlisting}
% \vspace*{-0.5cm} % Adjust for footer space
\caption{Obelisk persistence layer \textsc{ast}/\textsc{api} structures, from \texttt{/sur/server-state}.  (Continued from previous page.)}
\label{fig:obelisk_persistence_ast_2}
\end{figure}

The \lstinline[style=inlinecode]{database} structure reveals the simplicity of Obelisk time-travel. The database schema and content are both \lstinline[style=inlinecode]{+mop}s (ordered \lstinline[style=inlinecode]{+map}s) indexed by time. The magic of Urbit noun de-duplication does the rest.

\lstinline[style=inlinecode]{view-cache} provides for fast retrieval of view data after a view has been requested once. Any alteration of table data underlying a view will create a new entry keyed on namespace, view name, and time, but without content. The first request for that key will cause the entire \lstinline[style=inlinecode]{selection} to execute and save the content. There are currently only system views. User-defined views will be in a later release.

\lstinline[style=inlinecode]{table} defines the schema used to persist data for a persisted relation. \lstinline[style=inlinecode]{'column-lookup'} returns the column type (aura) as well as its position in the canonical ordering. \lstinline[style=inlinecode]{table}s require a unique primary index, and for now there are no secondary table indices.

\begin{lstlisting}[style=listingblock, label={lst:obelisk_table}]
+$  index
    $:  %index
        unique=?
        key=(list key-column)
==  ==
+$  key-column
    $:  %key-column
      name=@tas
      =aura
      ascending=?
==  ==
\end{lstlisting}
\noindent
\lstinline[style=inlinecode]{columns} returns the columns in canonical order (the order originally defined).

With the table definition established, the actual data is in a \lstinline[style=inlinecode]{file} structure, also containing relevant metadata. At the most basic level, every row of persisted data consists of a key---a list of atoms which is one or more of the subsequent columns, and the column data, which maps from column name to the atomic value. Of all the persistence layer structures, the one actually persisting data is probably the most prone to future development. For purposes of alternating query execution plans, data is persisted in more than one way. \lstinline[style=inlinecode]{(tree [(list @)} \lstinline[style=inlinecode]{(map @tas @)])} is an ordered map allowing ordered selection on partial keys. \lstinline[style=inlinecode]{column-catalog} is a currently abandoned attempt at persisting data column-wise. (Inserting large amounts of data stored column-wise proved impractical, so for now column-wise storage is deferred.) \lstinline[style=inlinecode]{indexed-rows} is the only form currently in use, so for the beta release every query plan execution is a table scan.

\subsection{Presentation Layer}

Aside from urQL, Obelisk users only ever interact with the presentation layer (Figure~\ref{fig:obelisk_presentation_ast}). The Obelisk app allows for two actions, \lstinline[style=inlinecode]{%tape}, which submits a command and default database name in a urQL \lstinline[style=inlinecode]{$tape}, and \lstinline[style=inlinecode]{%commands}, which is a list of \textsc{api} commands where every table and view is already fully qualified by database.

\begin{landscape}
\begin{figure}[p]
\centering
\vspace*{-0.8cm} % Adjust for header space
\begin{lstlisting}[style=listingblock, basicstyle=\setmonofont\footnotesize,numbers=none, xleftmargin=3em, multicols=3]
:: action
$%  $:  %tape
        default-database=@tas
        urql=tape
    ==
    $:  %commands
        cmds=(list command)
==  ==

:: cmd-result
$:  %results
    (list result)
==

:: vector
$:  %vector
    (lest vector-cell)
==

:: vector-cell
[p=@tas q=dime]
:: result
$%  [%message msg=@t]
    [%vector-count count=@ud]
    [%server-time date=@da]
    [%security-time date=@da]
    [%schema-time date=@da]
    [%data-time date=@da]
    [%result-set (list vector)]
==












:: command
$%  alter-index
    alter-namespace
    alter-table
    create-database
    create-index
    create-namespace
    create-table
    create-view
    delete
    drop-database
    drop-index
    drop-namespace
    drop-table
    drop-view
    grant
    revoke
    selection
    truncate-table
    update
==
\end{lstlisting}
% \vspace*{-0.5cm} % Adjust for footer space
\caption{Obelisk presentation layer \textsc{ast}/\textsc{api} structures, from \texttt{/sur/obelisk}.}
\label{fig:obelisk_presentation_ast}
\end{figure}
\end{landscape}

This section has covered most of the available commands; several planned commands like alters, drops, and security commands will parse but are not yet supported in the Obelisk engine.

Every urQL command returns a result structure of metadata and vectors, which we have already covered.

An untapped area of investigation is to change the \lstinline[style=inlinecode]{%result-set} from returning \lstinline[style=inlinecode]{(list vector)} to returning a \lstinline[style=inlinecode]{list} of trap cores which emit \lstinline[style=inlinecode]{vector} when evaluated. The lookup of columns row by row in the query relation is computationally intense and could be offloaded to the client side.

\section{Conclusion}

A database scripting language should provide the most intuitive and complete experience possible for organizing and querying data. Uniting the relational data model with set theory and predicate calculus has proven to be the most effective general means of organizing and retrieving data to date, however the hierarchical model still has its place and should be supported as far as possible. We have seen that even \textsc{sql} databases have at least a three-level hierarchy, four including namespaces, five when implemented as a temporal database.

Obelisk and its urQL scripting language represent a reimagining of the relational database paradigm, inspired by and tailored to the unique architecture of the Urbit computer while remaining adaptable to other computing environments and free of the constraints of \textsc{sql}’s backward compatibility, Obelisk leverages E.\ F.\ Codd’s foundational principles to address longstanding deficiencies in traditional \textsc{rdbms}s. Its innovations—time travel through database states, idempotent queries, proper set operations, elimination of nulls, and enforced composability—offer a streamlined, transparent, and powerful approach to data management. The integration of Urbit’s pure functional model ensures that all operations, whether through the urQL scripting language or the \textsc{api}, are consistent and accessible, with no hidden functionality. Additionally, through Urbit's global hierarchical namespace, data access via namespace path enhances flexibility, while Urbit’s identity-based security model provides robust protection. Obelisk not only advances the state of relational databases but also sets a new standard for clarity, consistency, and innovation in structured data systems.\tombstone

\section*{Appendix:  Examples}
\addcontentsline{toc}{section}{Appendix:  Examples}

\begin{lstlisting}[style=listingcode_sql]
/* The %sys database contains provenance metadata for
all the user databases.
*/

FROM sys.sys.databases
SELECT database, sys-agent, sys-tmsp, data-ship,
       data-agent, data-tmsp ;

/* Set the default database to animal-shelter for the
remaining queries. The rest of the schema metadata is
in %sys namespace views of each user database.
*/

FROM sys.tables
SELECT namespace, name, agent, tmsp, row-count ;

FROM sys.columns
SELECT namespace, name, col-ordinal, col-name, col-type ;

/* The %sys namespace views also tracks the schema and
content history.
*/

FROM sys.sys-log SELECT tmsp, agent, component, name ;

FROM sys.data-log
SELECT tmsp, ship, agent, namespace, table, row-count ;

/* Start exploring. Return complete sentences,
  if you like.
*/

FROM reference.calendar T1
JOIN reference.calendar-us-fed-holiday T2
WHERE T1.date BETWEEN ~2025.1.1
      AND ~2025.12.31
SELECT 'On this day: ', day-name, ',', T1.date,
       ', is the holiday ', us-federal-holiday, '.' ;

/* Create a new database, add a table and insert data.
The default namespace is %dbo.
*/

CREATE DATABASE db2 AS OF ~2000.1.1 ;

CREATE TABLE db2..my-table-1 (col1 @t, col2 @da)
PRIMARY KEY (col1) ;

INSERT INTO db2..my-table-1 (col1, col2)
  VALUES ('today', ~2000.1.1)
    ('tomorrow', ~2000.1.2)
    ('next-day', ~2000.1.3)
    ('next-today', ~2000.1.1)
    ('next-tomorrow', ~2000.1.2)
    ('next-next-day', ~2000.1.3) ;
\end{lstlisting}

\selectlanguage{USenglish}
\printbibliography
\end{document}
